\makeabstract
    {
    An Agent-Based Model of C. difficile Transmission Between Hospitals and Communities
    }
    {
    Hayden Curfman\textsuperscript{1}, Nora Hart\textsuperscript{2}, Natsuka Hayashida\textsuperscript{3}, Matthew Senese\textsuperscript{4}, Brittany Stephenson\textsuperscript{5}, Cara Sulyok\textsuperscript{6}
    }
    {
    \textsuperscript{1} Williams College, Williamstown, MA\\
\textsuperscript{2} Amherst College, Amherst, MA\\
\textsuperscript{3} Brown University, Providence, RI\\
\textsuperscript{4} University of Notre Dame, Notre Dame, IN\\
\textsuperscript{5} Lewis University, Romeoville, IL\\
\textsuperscript{6} Villanova University, Villanova, PA
    }
    {
    Clostridioides difficile (C. difficile) is a highly contagious bacteria that spreads through the fecal-oral route. C. difficile has typically been viewed as a nosocomial infection, but incidence has increased in the community in recent years. This model simulates the transmission of C. difficile between hospitals and communities to analyze the primary causes of colonizations in both the hospital and community, determine which mitigation strategies are the most effective in reducing incidence, and investigate the amount of time it takes for C. difficile to infiltrate the community from the hospital and vice versa.
    }
    {
    We developed an agent-based model using NetLogo that represents 10,000 people in a community including a hospital, houses, workplaces, schools, and public restrooms. Individuals can move between four disease classes: resistant, susceptible, colonized (asymptomatic), and diseased (symptomatic). To obtain baseline results, we ran 100 copies of one-year simulations in our community. Then, we implemented various mitigation strategies to see which most decreased incidence of colonization and disease.
    }
    {
    Preliminary results show that decreasing length of stay in the hospital is the most effective mitigation strategy for reducing incidence of colonization and disease in both the community and the hospital. The model shows that in the hospital, patients are more likely to become colonized with C. difficile through contact with contaminated surfaces rather than interactions with contaminated people. In the community, workplace interactions are the most common cause of colonization. Lastly, the mean time required for C. difficile to infiltrate the community from a contaminated hospital is approximately 12 times longer than the mean time required for the reverse.
    }
    {
    Our model highlights the potential impact of improving hospital efficiency on reducing the incidence of C. difficile not only in hospitals but also in communities. The prevalence of colonization associated with environmental surfaces in the hospital also highlights the importance of effective cleaning protocols and adherence to these protocols.
    }

    \newpage
    